% Appendix: the full fault table. Moved from Section~\ref{sec:conformance}. Rows found by "Kit" come
% from the campaign (evidence/conformance-campaign-run1.md, evidence/conformance-campaign-run2.md,
% investigations/conformance-run1-triage.md); the other rows come from those files. "Review" means
% the release notes attribute the fix to the project's own review or security audit, collectively
% rather than per item (2.0.0: "the fixes from the project's own security audit of every module";
% 1.4.0: "fail-closed gaps found by review"); "Accept." means the release notes attribute it to the
% acceptance scenarios (2.1.0). "n.s." means neither CHANGELOG.md (any version) nor docs/review.md
% attributes the row. Rows found by this paper's other experiments (Found by "\S n"): pgvector's scan
% limits from the Kaggle scale run, which the release notes call the acceptance run on Kaggle
% (CHANGELOG.md, the unreleased honest-completeness change); Chroma's crash from
% evidence/x3-x6-run1.md ("X3: Chroma, refusal and crash") and wardowl PR #187; SpiceDB staleness, the
% cache refill and the empty LookupResources stream from evidence/x4-x7-run1.md and x7-amendment2.md,
% PR #186; OPA's residual from evidence/x4-x7-run1.md (X4-H3); pgvector's plans from
% evidence/x8-run.md and investigations/x0-x1-pgvector-gap.md. PRs #167, #186 and #187 were merged on
% 2026-10-30; none is in a release yet (the latest is 2.1.0). The Upstream column was checked against
% vendor pages and trackers on 2026-10-30 (bib/STATUS.md).

\section{Faults and semantic gaps: the full table}\label{sec:conf-results}

Table~\ref{tab:conf-faults} lists the faults and semantic gaps in stores and engines that the
project's release notes, code comments and documentation record, up to the version we evaluate,
and those that this paper's experiments exposed. \emph{Found by} separates the two. Findings of this
paper's method are marked \emph{Kit}, the campaign of Appendix~\ref{sec:conf-campaign}, or \S$n$, the
section or appendix of another experiment. The project's history before this paper follows its release notes:
\emph{Review}, attributed collectively to the project's own code review or security audit;
\emph{Accept.}, to the acceptance scenarios against real servers; \emph{n.s.}, not stated. \emph{Effect} is the
consequence without the mitigation: \emph{Leak}, a false admit, which the re-check turns into a
violation; \emph{Leak*}, wrong stored or returned metadata, which also passes a re-check that reads
chunk metadata, because (A1) or (A2) fails; \emph{Leak\textdagger}, a check that itself admits,
beyond any re-check; \emph{Recall}, readable chunks silently missing; \emph{Stale}, a revoked grant
that still admits for a while after the change; \emph{Wrong rows}, a write reaching chunks it did
not name; \emph{Avail.}, requests failing closed; \emph{Perf.}, slower retrieval, as without a
pre-filter; and \emph{Crash}, the client process killed. PR~\#$n$ marks a mitigation merged on
2026-10-30 and not yet in a release. \emph{None (adapter)} marks a fault in \sys{}'s own adapter or
cache where the engine behaves as configured.

\begingroup
\scriptsize
\newcommand{\confrag}{\raggedright}
\setlength{\tabcolsep}{3pt}
\setlength{\LTcapwidth}{\linewidth}
\begin{longtable}{@{}p{0.115\linewidth}p{0.285\linewidth}p{0.075\linewidth}p{0.07\linewidth}p{0.225\linewidth}p{0.14\linewidth}@{}}
\caption{Faults and semantic gaps in vector stores (top) and authorization engines (bottom),
collected from the project's release notes, code comments and documentation and from this paper's
experiments; the legend is above. No row has been confirmed by the maintainers of the store or engine.}
\label{tab:conf-faults}\\
\toprule
Store or engine & Fault or semantic gap & Found by & Effect & Mitigation in \sys{} & Upstream \\
\midrule
\endfirsthead
\multicolumn{6}{@{}l}{\small Table~\ref{tab:conf-faults}, continued}\\
\toprule
Store or engine & Fault or semantic gap & Found by & Effect & Mitigation in \sys{} & Upstream \\
\midrule
\endhead
\midrule
\multicolumn{6}{r@{}}{continued on the next page}\\
\endfoot
\bottomrule
\endlastfoot
\confrag Elasticsearch, OpenSearch
 & \confrag Keyword fields compare as strings: a \texttt{terms} query for \texttt{1} matches
   \texttt{"1"}
 & \confrag n.s. & \confrag Leak
 & \confrag Type-tagged terms in one keyword field; metadata unindexed
 & \confrag Not stated~\cite{elastic2026keyword,opensearch2026keyword} \tabularnewline
\confrag pgvector (jsonb)
 & \confrag Numbers compare by value: $\mathrm{In}(\texttt{level},\{1\})$ matched a stored
   \texttt{1.0}
 & \confrag Review & \confrag Leak
 & \confrag Compare integers by the stored number's text
 & \confrag Documented~\cite{postgresql2026json} \tabularnewline
\confrag pgvector (jsonb)
 & \confrag A float of magnitude $\geq 10^{16}$, which Python's JSON encoder writes with an
   exponent, is stored and returned as an integer
 & \confrag Review & \confrag Leak*
 & \confrag Write such floats with a fraction
 & \confrag Documented~\cite{postgresql2026json} \tabularnewline
\confrag pgvector 0.8
 & \confrag Iterative HNSW scans stop at \texttt{hnsw.max\_scan\_tuples} or a memory limit: few or
   no rows under selective filters
 & \confrag \S\ref{sec:e2e} & \confrag Recall
 & \confrag Honest completeness; both limits documented
 & \confrag Limits documented~\cite{pgvector2026pgvector} \tabularnewline
\confrag pgvector via psycopg
 & \confrag Prepared statements: after five custom plans PostgreSQL may switch to a generic plan
   that reads the GIN index and sorts (exact, about 50 times slower) or keep custom HNSW plans
   (lossy), decided at run time
 & \confrag \S\ref{sec:qual-pgvector}, \S\ref{sec:e2e} & \confrag Recall or Perf.
 & \confrag None yet
 & \confrag Rule documented~\cite{postgresql2026prepare,psycopg2026prepared} \tabularnewline
\confrag Chroma 1.5.x, local mode
 & \confrag Array metadata of a deleted collection reappears on new records elsewhere and matches
   \texttt{where} filters (chroma-core/chroma\#7594)
 & \confrag n.s. & \confrag Leak*
 & \confrag No array metadata; metadata rebuilt from an exact JSON copy
 & \confrag Open; fix unmerged~\cite{sjmurdoch2026bug} \tabularnewline
\confrag Redis; Weaviate
 & \confrag Distinct values merge: TAG matching ignores case by default; \texttt{field}
   tokenization trims whitespace, joining \texttt{"hr "} and \texttt{"hr"}
 & \confrag n.s. & \confrag Leak
 & \confrag Hex-encoded terms (Redis); a separator at both ends of each term (Weaviate)
 & \confrag Redis: documented~\cite{redis2026tag}; Weaviate: in
   code~\cite{weaviate2026tokenizergo}, not docs~\cite{weaviate2026inverted} \tabularnewline
\confrag Redis 8.10
 & \confrag Tag values are cut at 4,096 bytes when indexed and looked up whole: long values never
   match, and a value equal to the kept prefix matches the longer one
 & \confrag Kit & \confrag Recall; Leak
 & \confrag Store and search a term over the limit as its SHA-256 digest (PR~\#167)
 & \confrag Reported 2026-10-30 (RediSearch issue \#12106) \tabularnewline
\confrag MongoDB Atlas 8.0, 8.3
 & \confrag Strings in vector-search filter fields are cut at 8,191 characters when indexed,
   undocumented (8,181 is documented for another field type): the same effects as for Redis
 & \confrag Kit & \confrag Recall; Leak
 & \confrag As for Redis (PR~\#167)
 & \confrag Documentation feedback sent 2026-10-30 \tabularnewline
\confrag Elasticsearch, OpenSearch, Milvus, Qdrant
 & \confrag Size and syntax limits refuse a write or query: 32,766-byte terms; array-element and JSON
   sizes; field-name grammar
 & \confrag Kit & \confrag Avail.
 & \confrag None needed: each fails closed with an error and returns nothing
 & \confrag Term limit~\cite{elastic2026ignore} and JSON size~\cite{milvus2026json} documented;
   Qdrant's grammar in its source \tabularnewline
\confrag Chroma 1.5.9, local mode
 & \confrag A flat \texttt{\$or} of about 1,000 clauses exceeds SQLite's expression depth; one of
   about 16,500 or more kills the client process (SIGSEGV)
 & \confrag Kit; \S\ref{sec:qual-size} & \confrag Avail.; Crash
 & \confrag Nest \texttt{\$and} and \texttt{\$or} 16 to a node; past four size bounds, filter 1,000
   candidates itself (PR~\#187)
 & \confrag Reported 2026-10-30 (Chroma issues \#7969, \#7970) \tabularnewline
\confrag Qdrant
 & \confrag \texttt{.}, \texttt{[}, \texttt{]} and \texttt{"} in a field name are path syntax:
   \texttt{team.name} reaches a nested object
 & \confrag n.s. & \confrag Leak
 & \confrag Refuse such field names
 & \confrag Path syntax documented~\cite{qdrant2026filtering} \tabularnewline
\confrag Pinecone, MongoDB, Milvus, LanceDB, Redis, Vespa
 & \confrag A missing or empty filter means ``match everything''
 & \confrag Review & \confrag Leak
 & \confrag Raise \texttt{NothingMatches}; send no query
 & \confrag Filters optional, e.g.~\cite{pinecone2026search,mongodb2026run} \tabularnewline
\confrag Elasticsearch, OpenSearch, Redis, Vespa
 & \confrag Partial results as success: failed or timed-out shards under HTTP 200; RESP3 timeout
   warnings; degraded coverage
 & \confrag Review & \confrag Recall
 & \confrag Disallow partial results; raise on each marker
 & \confrag Documented~\cite{elastic2026run,opensearch2026search,redis2026configuration,vespa2026graceful} \tabularnewline
\confrag Pinecone serverless; MongoDB Atlas
 & \confrag Writes become searchable only after a delay (eventual consistency; index lag)
 & \confrag n.s. & \confrag Stale
 & \confrag Writes wait until a search sees them (Pinecone: up to 30\,s)
 & \confrag Documented~\cite{pinecone2026check,mongodb2026configure} \tabularnewline
\confrag Cassandra 5.0
 & \confrag Paxos timestamps of lightweight transactions do not order with plain writes: an older
   permission could overwrite a newer one
 & \confrag Review & \confrag Leak*
 & \confrag Plain writes only
 & \confrag Documented by DataStax~\cite{datastax2026query} \tabularnewline
\confrag Milvus (pymilvus)
 & \confrag Reads and deletes by identifier quote identifiers into an expression: the identifier
   \texttt{a\textquotesingle,\textquotesingle b} also selects \texttt{a} and \texttt{b}
 & \confrag Review & \confrag Wrong rows
 & \confrag Identifiers as filter parameters
 & \confrag pymilvus PR \#3823 open~\cite{sungbin10152026fix} \tabularnewline
\midrule
\confrag OpenFGA
 & \confrag ListObjects stops at its result limit (1,000 by default) and returns a partial list
   without saying so
 & \confrag n.s. & \confrag Recall
 & \confrag At the limit, no filter: post-filter
 & \confrag Limits documented~\cite{openfga2026openfga} \tabularnewline
\confrag OpenFGA
 & \confrag ListObjects also stops at its deadline (3\,s by default), again without saying so
 & \confrag Review & \confrag Recall
 & \confrag After the deadline, no filter: post-filter
 & \confrag Limits documented~\cite{openfga2026openfga} \tabularnewline
\confrag SpiceDB
 & \confrag LookupResources may stop at the request limit (1,000 by default) or return results
   conditional on a caveat
 & \confrag n.s. & \confrag Recall; Leak
 & \confrag In either case, no filter: post-filter
 & \confrag Documented in the API~\cite{authzed2026permissionserviceproto} \tabularnewline
\confrag SpiceDB (adapter)
 & \confrag An empty LookupResources stream, which is also what a stream cut before its first line
   looks like, was read as ``nothing readable''
 & \confrag \S\ref{sec:eng-revoke} & \confrag Recall
 & \confrag No filter: post-filter, checking every candidate (PR~\#186)
 & \confrag None (adapter) \tabularnewline
\confrag SpiceDB (adapter)
 & \confrag Groups checked as separate subjects bypass exclusion: under
   \texttt{viewer - banned}, a banned user read through a group
 & \confrag Accept. & \confrag Leak\textdagger
 & \confrag Check the user only; SpiceDB resolves the memberships it stores
 & \confrag None (adapter) \tabularnewline
\confrag SpiceDB 1.56
 & \confrag Under \texttt{minimizeLatency}, a revoked permission admitted for 8.7\,s in 1 of 40
   trials, past the 5.5\,s its documented defaults imply, which we have not explained (a rerun's 40 trials: at most 5.2\,s), and
   could reappear after it had gone
 & \confrag \S\ref{sec:eng-revoke} & \confrag Stale
 & \confrag Declare 11\,s of staleness; after \texttt{invalidate()} the cache holds back that long
   (PR~\#186)
 & \confrag Bound documented~\cite{authzed2026consistency,authzed2026reference} \tabularnewline
\confrag SpiceDB, Cerbos (cache)
 & \confrag After \texttt{invalidate()}, \sys{}'s cache refilled from the engine's stale answer and
   held the revoked grant for its full ttl
 & \confrag \S\ref{sec:eng-revoke} & \confrag Stale
 & \confrag Cache nothing for the engine's declared staleness after \texttt{invalidate()}
   (PR~\#186)
 & \confrag None (adapter) \tabularnewline
\confrag OPA (partial evaluation)
 & \confrag The residual keeps \texttt{not (x in [])} even for users with no ban, and \sys{}'s
   filters have no negation: one ban rule leaves every user without a pre-filter
 & \confrag \S\ref{sec:eng-agree} & \confrag Perf.
 & \confrag None; such users are post-filtered
 & \confrag Not stated~\cite{opa2026evaluating,opa2026writing} \tabularnewline
\confrag Cerbos (adapter)
 & \confrag The translator maps CEL \texttt{in}, \texttt{hasIntersection} and \texttt{==} to In, which
   also matches a scalar and matches a list by element; in CEL, \texttt{in} or
   \texttt{hasIntersection} on a scalar is an error (denied) and \texttt{==} compares a list whole
 & \confrag Review; Kit & \confrag Leak
 & \confrag Translate \texttt{and}, \texttt{or}, \texttt{hasIntersection}, \texttt{in} and
   \texttt{==} only; else post-filter; documented
 & \confrag None (adapter); documentation note filed 2026-10-30 (Cerbos issue \#3456) \tabularnewline
\confrag Cedar (adapter)
 & \confrag Values Cedar cannot hold (floats, nulls, integers beyond 64 bits) and keys starting with
   \texttt{\_\_} were dropped, so a \texttt{forbid} on them stopped applying
 & \confrag Review & \confrag Leak\textdagger
 & \confrag Whole-number floats become Longs; a chunk with any other such value is denied
 & \confrag None (adapter) \tabularnewline
\confrag OpenFGA, SpiceDB
 & \confrag An identifier outside the engine's syntax makes it refuse the whole request
 & \confrag Review & \confrag Avail.
 & \confrag Deny such a chunk without sending it
 & \confrag Syntax in API specs~\cite{openfga2026openfgaproto,authzed2026coreproto} \tabularnewline
\confrag OpenFGA (adapter)
 & \confrag More than 100 contextual tuples, or a user or group identifier over its size limit, makes
   it refuse the whole request
 & \confrag Kit & \confrag Avail.
 & \confrag Refuse such a subject before sending, naming the limit (PR~\#167)
 & \confrag None (adapter) \tabularnewline
\confrag SpiceDB (adapter)
 & \confrag A subject or group identifier outside the object-identifier syntax makes it refuse the
   whole request
 & \confrag Kit & \confrag Avail.
 & \confrag Refuse such a subject before sending, naming the limit (PR~\#167)
 & \confrag None (adapter) \tabularnewline
\confrag Ory Keto (adapter)
 & \confrag Batches above \texttt{limit.max\_batch\_check\_size} (10 by default) are refused with
   HTTP 400; the adapter sent 100
 & \confrag Accept. & \confrag Avail.
 & \confrag Batches of 10
 & \confrag None (adapter) \tabularnewline
\end{longtable}
\endgroup

Of the table's 31 rows, 18 record the project's development before this paper and are history, not
findings of our method. Twelve are this paper's findings, six from the campaign and six from the
other experiments, and one, Cerbos, the campaign reproduced. Nine entries let a filter admit a chunk
the subject may not read, which leaks without
the re-check and becomes a violation with it, as Observation~\ref{prop:reduction} predicts. The five
Leak* and Leak\textdagger{} entries lie beyond any re-check: the kit's exact comparison of returned
with stored metadata targets the first kind, and only a test of the policy itself reaches the
second. The remaining effects cost recall, freshness, speed or availability, and one crashed the
client process; where a fix was needed, the lost recall had been reported as a complete answer, which
Appendix~\ref{sec:sys-complete} changes. Before the campaign, the record credited
review or the acceptance scenarios where it credited anything, never the kit.

Counted by kind, the campaign found two store behaviors with a correctness effect, Redis and MongoDB
cutting a long term at indexing, undocumented for the field type the layer uses; two availability
faults in the layer's own adapters, for OpenFGA and SpiceDB; and the documented limits of five stores
(Elasticsearch, OpenSearch, Milvus, Qdrant and Chroma), which fail closed. It reproduced one known
gap, in the Cerbos adapter's translator. All lie in what the generator targets, long values, many
groups, unusual identifiers and scalar fields, at sizes that the continuous-integration runs had not
reached. The filter-size experiment then showed that one of those limits, Chroma's, can also crash
the client, which is a store behavior of its own. The other experiments added SpiceDB's staleness
beyond its documented bound, in one unexplained trial, faults in \sys{}'s cache and SpiceDB adapter, a vacuous
negation that OPA's partial evaluation keeps, and pgvector's scan limits and its run-time choice
between a lossy and an exact plan. Of the eleven rows of Table~\ref{tab:conf-findings}, four are faults in \sys{}'s own adapters or cache. Several entries, the eight marked as
adapter faults among them, are faults in \sys{}'s own code that a store's or engine's behavior
exposed; any other integration meets the same behavior. Recall lost by approximate indexes under
selective or misaligned filters, such as Qdrant's filtered HNSW traversal, is measured in
Section~\ref{sec:quality} instead.
